\question
Для детерминированного конечного автомата задана таблица переходов (1 столбец – состояния, 1 строка – символы алфавита):

\begin{tabular}{||c | c |c| c||} 
 \hline
 & \textbf{a} & \textbf{b} & \textbf{c} \\ 
 \hline\hline
 \textbf{1} & 1& 1& 2\\ 
 \hline
 \textbf{2} & 3& 2& 1\\
 \hline
 \textbf{3} & 3& 4& 2\\
 \hline
 \textbf{4} & 3& 4& 4\\
 \hline
\end{tabular}

$S_0$ – 1, $F$ – {4}.

a) Постройте этот автомат в виде диаграммы Мура.

b) Автомату на вход поступила строка ‘bacccaab’. Определите, распознает ли автомат эту строку. Укажите, какие снапшоты (конфигурации) могут возникнуть в ходе обработки автоматом данного слова:

\begin{enumerate}
    \item ⟨2, ‘caab’⟩
    \item ⟨4, ‘b’⟩
    \item ⟨2, ‘ccaab’⟩
\end{enumerate}

---------------

Автор -- Тимур Гонтарь
